Kräfte auf die Leiter: ihre Gewichtskraft (in der Mitte), die Kraft $F_\mathrm{W}$ der Wand (horizontal, da die Wand glatt ist) und am Fuss die Normalkraft $F_\mathrm{N}$ und die Haftreibung $F_\mathrm{R}$ des Bodens. Sei $a$ der Abstand des Fusses von der Wand und $h$ die Höhe des oberen Endes.

Nimm den Fusspunkt als Drehachse: Die Kräfte des Bodens haben dann keinen Hebelarm. Die Wand drückt in der Höhe $h$ horizontal; die Gewichtskräfte haben als Hebelarm ihren horizontalen Abstand vom Fusspunkt.

Horizontal wirken nur die Wand und die Reibung, also muss die Reibung so gross sein wie die Kraft der Wand; vertikal trägt der Boden die ganze Gewichtskraft. Damit die Leiter nicht rutscht, braucht es $F_\mathrm{R}\le\mu\,F_\mathrm{N}$.

\begin{abcliste}
\abc
Die Gewichtskraft der Leiter hat den Hebelarm $a/2$:
\begin{align}
 F_\mathrm{W}\cdot h &= m\,g\cdot\frac{a}{2} \\
 F_\mathrm{W} &= \FWaF \\
   &= \frac{\mL\cdot\g\cdot\aF}{2\cdot\hW} \\
   &= \FWa \approx \result{\FWaP{0}{2}}
\end{align}
\begin{align}
 \mu &\ge \frac{F_\mathrm{R}}{F_\mathrm{N}} = \muaF \\
   &= \frac{\FWa}{\mL\cdot\g} \\
   &= \mua \approx \result{\muaP{0}{3}}
\end{align}

\abc
Die Person steht auf drei Vierteln der Leiter, also $f\,a$ ($f = \fPO$) vom Fusspunkt entfernt:
\begin{align}
 F_\mathrm{W}\cdot h &= m\,g\cdot\frac{a}{2} + M\,g\cdot f\,a \\
 F_\mathrm{W} &= \FWbF \\
   &= \frac{\left(\frac{\mL}{2}+\fP\cdot\MP\right)\cdot\g\cdot\aF}{\hW} \\
   &= \FWb \approx \result{\FWbP{0}{2}}
\end{align}
\begin{align}
 \mu &\ge \mubF \\
   &= \frac{\FWb}{(\mL+\MP)\cdot\g} \\
   &= \mub \approx \result{\mubP{0}{2}}
\end{align}
Je höher die Person steigt, desto mehr Reibung braucht die Leiter: Oben auf einer Leiter rutscht sie am ehesten. Auch eine flachere Leiter (grösseres $a$) braucht mehr Reibung.
\end{abcliste}
